\documentclass[11pt]{article}
\usepackage[margin=1in]{geometry}
\usepackage{hyperref}
\usepackage{enumitem}
\title{Responsible AI and Canva Publishing}
\author{Piter Garcia}
\date{\today}
\begin{document}
\maketitle

\section*{Grade and Class Match}
\textbf{Grade:} Grade 3\\
\textbf{Likely blocks:} Day 3 · 9:50-10:40 · Lamanaco; Day 4 · 9:50-10:40 · Lallucci; Day 5 · 9:50-10:40 · Regelsberger; Day 5 · 2:15-3:05 · Tandoi

\section*{Lesson Focus}
Responsible AI discussion paired with Canva-based publishing, design, and book-promotion work.

\section*{Learning Targets}
\begin{itemize}[leftmargin=*]
\item Students identify a responsible-use idea for AI-supported work.
\item Students create a visual product in Canva or a similar tool.
\item Students explain one design choice or message in their work.
\end{itemize}

\section*{Materials}
\begin{itemize}[leftmargin=*]
\item Student devices
\item Canva or equivalent tool
\item Prompt or example visuals
\end{itemize}

\section*{Suggested Lesson Flow}
\begin{enumerate}[leftmargin=*]
\item Open with the exact device, tool, or routine students need in order to start successfully.
\item Model one example task before students begin independent or partner work.
\item Give students time to build, code, design, or document their work while circulating for support.
\item Close with a short share-out, reflection, or debugging discussion.
\end{enumerate}

\section*{Source Notes}
This lesson packet is enriched with transcript-derived lesson notes, the school schedule roster, and official starting-point resources. See the paired Markdown guide for clickable links.

\bibliographystyle{plain}
\bibliography{responsible-ai-and-canva-publishing}
\end{document}
